	\noindent УДК 004.932 

\vspace{0.5em}
\noindent \textbf{Липянчик Євген} \\
Державний університет «Київський авіаційний інститут», Київ \\
Науковий керівник – \textbf{Василик Віталій, доктор фізико-математичних наук}

\vspace{1em}
\begin{center}
	\textbf{ДОСЛIДЖЕННЯ ТА ПРОГРАМНА РЕАЛIЗАЦIЯ МЕТОДУ МАСШТАБУВАННЯ
		ЗОБРАЖЕНЬ З ВИКОРИСТАННЯМ БАРИЦЕНТРИЧНОЇ ФОРМУЛИ НА ОСНОВI
		ВУЗЛIВ IНТЕРПОЛЯЦIЇ ЧЕБИШЕВА}
\end{center}

\vspace{0.5em}
\noindent \textbf{АНОТАЦІЯ.} Дослідження присвячене порівняльному аналізу ефективності алгоритмів масштабування зображень. Проаналізовано класичні локальні методи (Nearest Neighbor, Bilinear, Bicubic, Lanczos4) та метод інтерполяції за допомогою барицентричної формули на основі вузлів Чебишова [2]. Розроблено програмне забезпечення на базі C++ [3] та фреймворку Qt [4] для візуалізації та вимірювання часу виконання алгоритмів.

\noindent \textbf{Ключові слова:} інтерполяція, масштабування зображень, друга барицентрична формула, вузли Чебишова, OpenCV, багатопотоковість.

\textbf{Вступ.} Задачі масштабування зображень є критично важливими у сферах геоінформаційних систем (ГІС), обробки супутникових знімків, медичної візуалізації та комп'ютерного зору. Від вибору алгоритму масштабування безпосередньо залежить баланс між обчислювальною складністю та якістю зображення [2]. Інтерполяційні формули за вузлами Чебишева широко використовуються в математиці (див. на пр. [5][6]), і також можуть бути використані для задачі масштабування зображень.

\textbf{Мета.} Порівняти обчислювальну ефективність та візуальну якість барицентричної інтерполяції зі стандартними алгоритмами бібліотеки OpenCV [1] під час масштабування зображень.

\textbf{Матеріали та методи.} Для проведення експериментів розроблено програмне забезпечення мовою C++ [3] із використанням графічного фреймворку Qt [4]. Як еталонні алгоритми використано методи бібліотеки OpenCV [1] (cv::INTER\_NEAREST, \allowbreak cv::INTER\_LINEAR, \allowbreak cv::INTER\_CUBIC, \allowbreak cv::INTER\_LANCZOS4). Як альтернативу було реалізовано метод барицентричної інтерполяції з використанням вузлів Чебишова [2]. Для компенсації високої обчислювальної складності глобального методу розроблено оптимізовану архітектуру. Тестування продуктивності проводилося на різних коефіцієнтах масштабування.

\textbf{Математична модель.}
У даному розділі розглядається математичний опис задачі масштабування. З теоретичної точки зору, зображення $\mathcal{I}$ можна представити як двовимірну функцію $f(x,y)$, аргументи якої належать квадрату $A = [-1, 1]^2$ [2]. Цифрове представлення зображення $I$ розміром $n \times m$ пікселів розглядається як набір значень цієї функції, обчислених на дискретній сітці вузлів $X_{n \times m} \subset A$. Вважається, що дана сітка формується як тензорний добуток одновимірних систем вузлів: $X_{n \times m} = X_n \times X_m$, де кожна система визначається як:
\begin{equation}
	X_{\mu} := \{ x_k^\mu : k = 0, \dots, \mu \} \subset [-1, 1], \quad \forall \mu \in \mathbb{N}.
\end{equation}

У цій роботі для дискретизації використовуються вузли Чебишева II роду:
\begin{equation}
	x_i = -\cos\left(\frac{i\pi}{n}\right), \quad i = 0, \dots, n.
\end{equation}

Ці точки визначають сітку $X_{n \times m} = \{ (x_i^n, y_j^m) : i=0:n, j=0:m \}$, на якій задано вихідне зображення $I = f(X_{n \times m})$. У реальних умовах вхідні дані $I^{in}$ можуть бути дещо спотвореною версією ідеальної функції:
\begin{equation}
	I_{i,j}^{in} := \tilde{f}(x_i^n, y_j^m), \quad i = 0, \dots, n, \quad j = 0, \dots, m.
\end{equation}

Задача масштабування полягає у знаходженні апроксимації значень функції на новій сітці $X_{N \times M}$ (більш щільній для збільшення або розрідженій для зменшення):
\begin{equation}
	I_{k,h}^{res} := f(x_k^N, y_h^M), \quad k = 0, \dots, N, \quad h = 0, \dots, M.
\end{equation}

На відміну від роботи [2] ми використовуємо математичну модель по вузлах Чебишова 2-го роду. Для ефективної реалізації інтерполяції ми узагальнили одновимірну другу барицентричну формулу на випадок  2D:
\begin{equation*}
	L_{n,m}(x_r, y_r) = \frac{\sum\limits_{i=0}^{n}{}^{\prime} \sum\limits_{j=0}^{m}{}^{\prime} \frac{(-1)^{i+j}}{(x_r-x_i)(y_r-y_j)} f_{ij}}{\sum\limits_{i=0}^{n}{}^{\prime} \sum\limits_{j=0}^{m}{}^{\prime} \frac{(-1)^{i+j}}{(x_r-x_i)(y_r-y_j)}}
\end{equation*}
Символ $\sum'$ означає:
\begin{equation*}
	\sum_{p=0}^{N}{}^{\prime} a_p = \frac{1}{2} a_0 + \sum_{p=1}^{N-1} a_p + \frac{1}{2} a_N
\end{equation*}

\textbf{Висновки.} 
У результаті виконання роботи було розроблено та програмно реалізовано новий метод для масштабування цифрових зображень, що базується на використанні другої інтерполяційної барицентричної формули за вузлами Чебишова [2]. Проведено комплексне порівняння запропонованого алгоритму з існуючими базовими методами бібліотеки OpenCV [1].
Результати числових експериментів засвідчили, що при збільшенні зображень (upscaling) запропонований метод демонструє якість, порівнянну з класичною бікубічною інтерполяцією. Водночас, при зменшенні зображень (downscaling) алгоритм забезпечує значно кращі показники PSNR та SSIM [2].

\newpage
\noindent \textbf{Список використаних джерел:}
\begin{enumerate}
	\setlength{\itemsep}{0pt}
	\item Документація бібліотеки OpenCV (Image Geometric Transformations). URL: \url{https://docs.opencv.org/4.x/da/d54/group__imgproc__transform.html}.
	\item Occorsio D., Ramella G., Themistoclakis W. Lagrange-Chebyshev Interpolation for image resizing // Mathematics and Computers in Simulation. 2022. Vol. 197. P. 105–126.
	\item Bjarne Stroustrup "Programming principles and practice using C++".
	\item Qt Documentation. URL: \url{https://doc.qt.io/}.
	\item Gavrilyuk I.P., Makarov V.L., Vasylyk V.B. Exponentially convergent algorithms for abstract differential equations. Birkhauser, 2011.
	\item V. Vasylyk. Exponentially convergent method for differential equation in Banach space... // Journal of Computational \& Applied Mathematics. 2016.
\end{enumerate}
